%% Research Note of the AAS (K2 single events x TESS), AASTeX v7.0.1.
%% Source, numbers and provenance: https://github.com/coracle-us/duo-transits (k2singles/rnaas/).
\documentclass{aastex701}

\begin{document}

\title{Orbital Periods for K2 Single-transit and Single-eclipse Events from K2 Revisits and TESS}

\author[orcid=0009-0000-7524-4470]{Jon McBee}
\affiliation{Independent researcher, USA}
\email{jon.mcbee@gmail.com}

\begin{abstract}
\citet{LaCourse2018} catalogued 164 K2 events seen only once. I searched all later K2 campaigns
and all TESS data for repeats of the 144 events that TESS should detect. Every period I report
comes from fitted eclipse times and was re-derived blind with independent code. I link K2 C5
and C16 transits to the TESS transits of TOI-5560.01, narrowing its period from $\leq702$\,d to
$140.4939$ or $70.2470$\,d. The K2 planet candidate EPIC~211503363 has $P=122.905$ or
$61.453$\,d, but its companion is $\approx2.2\,R_{\rm J}$, probably an M dwarf. I give new or
refined periods or short alias lists for nine eclipsing binaries. These include HD~220967, a
bright, eccentric ($e\geq0.58$), 73.245-d system. In another, EPIC~202126877, the primary eclipse
deepened from 6\% in 2014 to 21\% in 2021 and 36\% in early 2026.
\end{abstract}

\keywords{\uat{Exoplanet detection methods}{489} --- \uat{Transit photometry}{1709} --- \uat{Eclipsing binary stars}{444}}

\section{Sample and Methods}

The catalogue lists 164 events on 163 stars; 158 fall in TESS sectors already observed, 145 of
those have an expected TESS signal-to-noise ratio of at least 7 per event (from the catalogued depth
and duration), and 144 have a physical depth. I combined K2
PDCSAP light curves \citep{Howell2014} with TESS SPOC, TESS-SPOC \citep{Caldwell2020} and QLP
\citep{Huang2020} light curves, plus my own TESScut photometry \citep{Brasseur2019} where pipeline
products are missing. Against SPOC data for ten stars, it gives depths within 3\% for uncrowded
stars (about 30\% shallower for crowded ones) and mid-times within 2\,min.

A box of the catalogued depth and duration slides along the data; a repeat needs SNR $\geq8$,
0.5--2 times the catalogued depth and 80\% coverage. The matching is deliberate, since only a
recurrence of the same event constrains $P=\Delta T/n$; it finds secondaries of similar depth but
misses binaries with only a shallow secondary. A period is rejected if any predicted epoch falls on
flat data. For eccentric binaries, primaries and secondaries were assigned jointly.

Final periods come from trapezoid fits to each eclipse and a linear ephemeris. Three blind analyses
with independent code, starting from the catalogue epoch alone, re-derived every period in
Table~\ref{tab:k2}; apart from EPIC~202126877, which shows timing variations, they agree within
$2\sigma$. Sanders~617 in M67 comes out at $62.59482\pm0.00003$\,d, against $62.59484\pm0.00003$\,d
from \citet{Sandquist2021}, with its C5 eclipse timed in EVEREST because PDCSAP removes its core. I
checked each star against ExoFOP, VizieR, SIMBAD, the Villanova K2 EB catalogue and Gaia DR3
\citep{Gaia2023,Arenou2023}.

\section{Results}

The catalogued event is present at its epoch in K2 PDCSAP (a dip deeper than half the catalogued
depth at SNR $>5$) for 123 of the 144 events. Of the other 21, checked in SAP, EVEREST, K2SFF and
K2SC, PDC removed or distorted eight, usually by blanking the eclipse core; one catalogued epoch is
a secondary 7.6\,d from the primary; four are shallower or weaker than required; one is hidden by
roll systematics; and seven appear only in K2SFF, with aperture-dependent depths, so they are
probably artefacts or blends.

Of the 144 events, 62 show at least one candidate repeat and 16 have a period that explains two or
more; 14 are in Table~\ref{tab:k2}. EPIC~202137580 is already well catalogued, the giant
EPIC~212694013 keeps over ten aliases, and HD~220967 needed a joint ephemeris. Events with a single
repeat were not vetted systematically.

\textit{TOI-5560.01 (EPIC~211924561).} Transits in TESS Sectors 46 and 72 left
$P\leq702.47$\,d. The K2 C5 transit and a C16 transit, masked in EVEREST
and K2SFF but clear in PDCSAP and SAP, fit one linear ephemeris with them. The cycle counts (7, 17,
22) share no common factor, and flat data exclude every shorter period except half:
$P=140.49392\pm0.00014$ or $70.24696\pm0.00007$\,d, $T_0={\rm BJD_{TDB}}\,2457179.8979\pm0.0014$.
The V-shaped transit (2.3\,ppt, 3.0\,h) is far shorter than a central transit of this F star
(12--15\,h), so the orbit is grazing, eccentric or both. Gaia DR3 radial velocities are not variable
($p=0.08$), implying $M_2\sin i\lesssim0.1\,M_\odot$ for a circular orbit. Both periods predict a
transit at ${\rm BJD_{TDB}}\,2461394.716\pm0.004$ (2026 December 20, 05:03\,UTC); one at 2461464.963
(2027 February 28) would confirm 70.25\,d.

\textit{EPIC~211503363.} This single-transit candidate of \citet{Kovacs2020} transits again in
Sectors 45 and 72 (3.4--3.8\,ppt, 23\,h), leaving
$P=122.9052\pm0.0006$ or $61.4526\pm0.0003$\,d. The host is evolved ($R_\star=3.6\,R_\odot$,
$\log g=3.46$; Gaia DR3), so the companion is $\approx2.1$--$2.4\,R_{\rm J}$: too large for a planet
or brown dwarf on a non-irradiated orbit, and consistent with an M dwarf. The $R_p=0.260$ of
\citet{Kovacs2020}, in a column labelled Jupiter radii, is in solar radii
($\sqrt{0.004019}\times4.098\,R_\odot$).

\textit{Eclipsing binaries.} HD~220967 ($G=9.4$, F7/8\,V) shows eclipses 37\% and 30\% deep in K2
C12 and TESS. The secondary at phase 0.867 gives $e\geq0.58$, and its 1.9 times longer duration
suggests $e\approx0.66$. The eclipses of the red
giant EPIC~211892898, noted by \citet{Libralato2016}, fall at the primary conjunctions of its Gaia
DR3 SB1 orbit ($72.53\pm0.03$\,d) within $1.8\sigma$. For EPIC~211489484 my ephemeris excludes a
106.94-d SPOC TCE \citep{Valizadegan2025}.

EPIC~202126877 is a known 21.78-d EB \citep{IJspeert2024}. Its K2 eclipse lies on cycle $-124$ of
the 2021--2023 TESS ephemeris but 2.2\,h early, a $3\sigma$ offset given the 40-min propagated
uncertainty. It is only 6.2\% deep in EVEREST and K2SFF, with no secondary ($<0.1\%$). In my
fixed-aperture TESS photometry (flux within $\pm15$\,min of mid-eclipse) the primary is 21\% deep in
2021, 27\% in 2023 and 36\% in January 2026, and the secondary grows from 4\% to 12\% between 2021
and 2023. TESS-SPOC (2021) and QLP (2021, 2023) agree within 5\%; the 2026 depth rests on my
photometry and agrees within 1\% across four apertures. The secondary-to-primary ratio doubles from
0.21 to 0.46, which a constant dilution cannot do, and the K2 aperture is uncrowded (CROWDSAP
0.999). A precessing orbit, perhaps driven by a third body, could cause both changes. The dedicated
fits and blind analyses, which use no depth window, recovered every eclipse.

For the C0 binaries, my 34.7456-d period for EPIC~202072917 agrees with \citet{IJspeert2024} and a
SPOC TCE \citep{Valizadegan2025}, and my 22.9041-d period for EPIC~202071902 agrees with
\citet{Guo2025} and is $3\sigma$ from the SB2 period of \citet{Kovalev2024}. Their VSX periods (3.842
and 8.697\,d) would put several eclipses in every TESS sector.

\begin{deluxetable*}{lllll}
\tabletypesize{\scriptsize}
\tablecaption{Periods from K2 revisits and TESS.\label{tab:k2}}
\tablehead{\colhead{EPIC} & \colhead{Events$^e$} & \colhead{Period (d)} & \colhead{Prior period (d)} & \colhead{Notes}}
\startdata
\cutinhead{Planet candidates (as catalogued)}
211924561 & C5, C16; S46, S72 & 140.49392(14), 70.246961(70) & $\leq702.47$ (TOI) & TOI-5560.01 \\
211503363 & C5; S45, S72 & 122.90522(64), 61.45261(32) & none$^g$ &  \\
\cutinhead{Validation (published orbit)}
211411112 & C5, C16; S46 & 62.594821(27) & 62.59484(3)$^a$ & Sanders 617 (M67) \\
\cutinhead{Eclipsing binaries: new or refined periods}
246194159 & C12; S70, S92 & 73.2451778(74) & none & HD 220967; sec.\ 0.867 \\
211489484 & C5, C18; S7, S46 & 82.258438(52) & TCE 106.94 (excl.) & sec.\ 0.194 \\
211892898 & C5, C16, C18 & 72.5089(14) & Gaia SB1 72.53(3) &  \\
248811085 & C14; S46, S72 & 49.250131(13) & none &  \\
246089124 & C12; S70, S92 & 73.962413(20) & none & sec.\ 0.494 \\
246696483 & C13; S43, S70 & 182.24526(18), 91.122630(90) & none & $P/2$ disfavoured \\
211390677 & C5, C18; S45 & 184.3109(11), 92.15543(55) & none &  \\
247967714 & C13; S44, S71 & 185.626676(21)/$n$, $n=1$--$3$ & none &  \\
212012030 & C5, C18; S71 & 220.321779(57)/$n$, $n=1$--$5$ & none &  \\
\cutinhead{Eclipsing binaries: published periods}
202126877 & C0; S43--45, S71--72 & 21.775 & 21.78$^b$ & deepening$^f$ \\
202072917 & C0; S6, S43, S45 & 34.7455865(70) & 34.744$^{b}$, 34.7457$^{c}$ &  \\
202071902 & C0; S44, S71 & 22.9041282(30) & 22.859(15)$^d$ &  \\
\enddata
\tablecomments{Uncertainties in the last digits in parentheses.
Without a secondary eclipse (``sec.'', at the given phase), twice the period cannot be excluded.
$^a$\citet{Sandquist2021}; $^b$\citet{IJspeert2024}; $^c$SPOC TCE \citep{Valizadegan2025};
$^d$\citet{Kovalev2024}; see also \citet{Guo2025}. $^e$K2 campaigns (C) and TESS sectors (S) fitted.
$^f$Primary 6\% deep in 2014, 21--36\% in 2021--2026. $^g$\citet{Kovacs2020}.}
\end{deluxetable*}

\section{Conclusion}

Most K2 single events with repeats are eclipsing binaries, but the search also narrows a TESS planet
candidate to two periods and shows that a K2 planet candidate is probably an M dwarf. Code and fits are at \url{https://github.com/coracle-us/duo-transits}.

\begin{acknowledgments}
This work uses data from MAST: K2 \dataset[doi:10.17909/t9ws3r]{https://doi.org/10.17909/t9ws3r},
SPOC \dataset[doi:10.17909/t9-nmc8-f686]{https://doi.org/10.17909/t9-nmc8-f686}, TESS-SPOC
\dataset[doi:10.17909/t9-wpz1-8s54]{https://doi.org/10.17909/t9-wpz1-8s54} and QLP
\dataset[doi:10.17909/t9-r086-e880]{https://doi.org/10.17909/t9-r086-e880} light curves and TESS
full-frame images \dataset[doi:10.17909/0cp4-2j79]{https://doi.org/10.17909/0cp4-2j79}; Gaia DR3;
ExoFOP \dataset[doi:10.26134/ExoFOP5]{https://doi.org/10.26134/ExoFOP5}; VizieR; SIMBAD; and the
Villanova K2 EB catalogue.
\end{acknowledgments}

\facilities{Kepler (K2), TESS, Gaia}
\software{lightkurve \citep{Lightkurve2018}, astroquery \citep{Ginsburg2019}, astrocut \citep{Brasseur2019}}

\section*{AI disclosure statement}
I used Anthropic's Claude, an AI assistant, under my supervision to write and run the analysis code
(including blind re-analyses and adversarial reviews by separate AI agents) and to draft this note. I
checked every number against the saved outputs and every reference against ADS, arXiv, Crossref or
VizieR, and take full responsibility for the content.

\begin{thebibliography}{}
\bibitem[Brasseur et al.(2019)]{Brasseur2019} Brasseur, C.~E., Phillip, C., Fleming, S.~W., Mullally, S.~E., \& White, R.~L. 2019, Astrocut: Tools for creating cutouts of TESS images, Astrophysics Source Code Library, ascl:1905.007
\bibitem[Caldwell et al.(2020)]{Caldwell2020} Caldwell, D.~A., Tenenbaum, P., Twicken, J.~D., et al. 2020, RNAAS, 4, 201
\bibitem[Gaia Collaboration, Arenou et al.(2023)]{Arenou2023} Gaia Collaboration, Arenou, F., Babusiaux, C., et al. 2023, A\&A, 674, A34
\bibitem[Gaia Collaboration, Vallenari et al.(2023)]{Gaia2023} Gaia Collaboration, Vallenari, A., Brown, A.~G.~A., et al. 2023, A\&A, 674, A1
\bibitem[Ginsburg et al.(2019)]{Ginsburg2019} Ginsburg, A., Sip{\H o}cz, B.~M., Brasseur, C.~E., et al. 2019, AJ, 157, 98
\bibitem[Guo et al.(2025)]{Guo2025} Guo, S., Kovalev, M., Li, J., et al. 2025, ApJS, 278, 46
\bibitem[Howell et al.(2014)]{Howell2014} Howell, S.~B., Sobeck, C., Haas, M., et al. 2014, PASP, 126, 398
\bibitem[Huang et al.(2020)]{Huang2020} Huang, C.~X., Vanderburg, A., P\'al, A., et al. 2020, RNAAS, 4, 204
\bibitem[IJspeert et al.(2024)]{IJspeert2024} IJspeert, L.~W., Tkachenko, A., Johnston, C., \& Aerts, C. 2024, A\&A, 691, A242
\bibitem[Kov\'acs(2020)]{Kovacs2020} Kov\'acs, G. 2020, A\&A, 643, A169
\bibitem[Kovalev et al.(2024)]{Kovalev2024} Kovalev, M., Zhou, Z., Chen, X., \& Han, Z. 2024, MNRAS, 527, 521
\bibitem[LaCourse \& Jacobs(2018)]{LaCourse2018} LaCourse, D.~M., \& Jacobs, T.~L. 2018, RNAAS, 2, 28
\bibitem[Libralato et al.(2016)]{Libralato2016} Libralato, M., Nardiello, D., Bedin, L.~R., et al. 2016, MNRAS, 463, 1780
\bibitem[Lightkurve Collaboration et al.(2018)]{Lightkurve2018} Lightkurve Collaboration, Cardoso, J.~V.~d.~M., Hedges, C., et al. 2018, Lightkurve: Kepler and TESS time series analysis in Python, Astrophysics Source Code Library, ascl:1812.013
\bibitem[Sandquist et al.(2021)]{Sandquist2021} Sandquist, E.~L., Latham, D.~W., Mathieu, R.~D., et al. 2021, AJ, 161, 59
\bibitem[Valizadegan et al.(2025)]{Valizadegan2025} Valizadegan, H., Martinho, M.~J.~S., Jenkins, J.~M., et al. 2025, AJ, 170, 287
\end{thebibliography}

\end{document}
